% Section 6.4, from lectures/L06/appendix/d_what_to_build.md.

\section{What to Build}
\label{c6:sec:build}

\subsection{The task}
\label{c6:sec:task}
Two things, and the second of them is the book's capstone: \file{drivers/source/watchdog.asm},
which writes the watchdog's control byte correctly; and a C program that calls the assembly you
wrote in Chapter~1, unchanged.

\subsection{What to build in assembly: \code{drivers/source/watchdog.asm}}
\label{c6:sec:driver}
Four subroutines.

\textbf{\code{watchdog_init}} takes the control byte in \code{r24} and performs the timed sequence
(\secref{c6:sec:timed}):
\begin{enumerate}
  \item Save \code{SREG}, then \code{cli}.
  \item \code{wdr}, so the timeout starts from now rather than from wherever the counter had got to.
  \item Clear \code{WDRF} in \code{MCUSR}, leaving the other three cause bits alone.
  \item Write \code{WDCE} and \code{WDE} together into \code{WDTCSR}.
  \item Write \code{r24} into \code{WDTCSR}.
  \item Restore \code{SREG}.
\end{enumerate}

\textbf{\code{watchdog_reset}} is a single \code{wdr} and a \code{ret}. It runs in your main loop,
thousands of times, and must do exactly that one thing.

\textbf{\code{watchdog_disable}} is \code{watchdog_init} with a written value of zero.

\textbf{\code{watchdog_caused_reset}} returns 1 in \code{r24} if \code{WDRF} was set in
\code{MCUSR}, else 0, and clears \code{WDRF} either way, leaving the other cause bits alone.

\subsubsection{Details that are pinned rather than up to you}
\textbf{Steps 4 and 5 must be adjacent.} Two \code{sts} instructions, nothing between them. Two
cycles each, so they fit the four-cycle window with one cycle to spare; one single-cycle
instruction between them still fits, and anything more leaves \code{WDTCSR} in system reset mode,
holding \code{0x08} on a device fresh from reset (\secref{c6:sec:timed}).

\textbf{Save and restore \code{SREG}, do not just \code{cli} and \code{sei}.} The sequence must run
with interrupts off, but leaving them off on the way out is a bug: every other interrupt in the
program quietly stops, at a moment determined by whoever last configured the watchdog. This is the
same discipline as Chapter~3's atomic read (\secref{c3:sec:shared}).

\textbf{Clear \code{WDRF} before the window, not after.} The hardware refuses to clear \code{WDE}
while \code{WDRF} is set, so a driver that clears it afterwards works the first time and cannot
disable the watchdog after a watchdog reset, which is exactly when it matters.

\textbf{\code{MCUSR} is inside the I/O window, so \code{in} and \code{out} reach it}; \code{WDTCSR}
is not, and needs \code{lds} and \code{sts}. The two numbers are the usual pair: \code{MCUSR} is
data space \code{0x54} and I/O \code{0x34}, and the name in \file{m328Pdef.inc} is the I/O one, so
\code{in r24, MCUSR} is what you write (\secref{c1:sec:iowindow}). \code{WDTCSR} is \code{0x60} in
both, because above the window there is only one. Two registers, two addressings, one subroutine.

\textbf{Leave the other reset causes alone.} \code{MCUSR} records four different reasons in four
bits, and clearing all of them because you cared about one destroys information nothing can
recover.

\subsection{The capstone: call it from C}
\label{c6:sec:capstone}
Write \file{drivers/app/main.c}, a C program that calls \code{shift_bits} from Chapter~1 and uses
the result.

\textbf{Change nothing in \file{utils.asm}.} That is the point. It has followed the calling
convention since Chapter~1, so a C prototype is the whole of the interface
(\secref{c6:sec:asmfromc}).

Build it with a normal \code{avr-gcc} invocation rather than the book's usual one: a C program
needs the startup code that sets the stack pointer, zeroes \code{.bss} and \code{r1}, and calls
\code{main}, all the things your \file{main.asm} did by hand.

Then generate the compiler's own version of the same job with \code{-S}, and compare it against
what you wrote (the cross-check, Exercise~\ref{c6:ex:crosscheck}).

\subsubsection{Details that are pinned rather than up to you}
\textbf{The prototype's types must match.} \code{uint8_t shift_bits(uint8_t)}. Declaring the
argument wider passes it in \code{r25:r24} and your subroutine reads only \code{r24}, which changes
nothing you can see and is still wrong; declaring the result wider as well reads back an
\code{r25} your subroutine never wrote (\secref{c6:sec:asmfromc}).

\textbf{Do not add \code{extern "C"}.} This is C, not C++, and there is no name mangling to
suppress.

\subsection{What none of this does}
\label{c6:sec:notdo}
\textbf{The driver does not choose the byte.} It writes whatever \code{r24} holds. Working that
byte out from \secref{c6:sec:registers} is your program's job, and so is deciding to write it at
all.

\textbf{The driver does not decide when to pet.} Where \code{watchdog_reset} is called from is the
entire difference between a watchdog that works and one that is decoration
(\secref{c6:sec:detects}), and nothing in the driver can help.

\textbf{And none of it is tested against sleep.} The interaction between sleep modes and wake
sources is specified by the datasheet and approximated by the simulator, so this book teaches it
and does not check it (\secref{c6:sec:sleepnot}).
